% ================================================================
\chapter{Introduction}
% ================================================================
\begin{readaloud}
Every table has its own story. What follows is a system for
telling it --- one set of rules that travels as far as your
imagination will carry it.
\end{readaloud}

\lettrine{L}{egend of the Technomancers} is a tabletop roleplaying
game about travellers who move between worlds. The setting
shifts --- high fantasy one night, rain-slicked cyberpunk the next,
the hard vacuum of deep space the one after that --- but the
characters persist. Their scars, their debts, their hard-won skills:
all of it comes with them. The rules are simple enough to learn
in one session and flexible enough to carry a campaign of a hundred.

One player is the \keyword{Game Master} (GM), who builds the world,
voices its inhabitants, and decides what the dice mean. Everyone
else plays a \keyword{Traveller} --- an exceptional individual
navigating whatever the GM puts in front of them. Travellers are
not superheroes. They are people who have paid attention,
practised hard, and made the choice to step through the wound in
the world when it opened.

\section{What This Game Is About}

At the heart of \emph{Legend of the Technomancers} is a single
question: \emph{what kind of person are you when the world changes?}
The mechanics are built around that question. Your character's
Stats --- \keyword{Body}, \keyword{Mind}, \keyword{Face} --- do not
change when you cross a threshold into a new genre. Your skills
travel with you. Your magic tags re-dress themselves in the
vocabulary of wherever you have landed, but the power behind them
is the same power.

The game rewards players who think about their characters as
\emph{people} first and as stat blocks second. The rules will handle
the numbers. Your job is to decide what your Traveller wants, fears,
and would sacrifice everything for.

\section{How This Book Is Organised}

\begin{itemize}
  \item \keyword{Part I: Building a Traveller} covers character
        creation and a chapter for each of the four columns ---
        Body, Stuff, Skills, and Magic.
  \item \keyword{Part II: Playing the Game} covers the core dice
        engine in full, then the shape of a session: turns, scenes,
        GM tools, and advancement.
  \item \keyword{Part III: The Realms} presents three example
        settings --- High Fantasy, Cyberpunk, and Space Adventure
        --- and shows how the same rules express themselves in
        each one.
\end{itemize}

If you want the short version of how rolls work before diving into
character creation, the core engine is explained in full at the
start of Part~II. The rules in Part~I give you everything you need
to build your Traveller; Part~II explains what to do with them at
the table.

\section{What You Need to Play}

You need a handful of six-sided dice --- a dozen is comfortable,
twenty is luxurious. Every player needs a character sheet (or a
notebook that can serve as one). The GM needs a screen or a
notebook, something to track initiative and hit points behind, and
whatever maps or reference cards feel useful.

Everything else --- imagination, willingness to fail in interesting
ways, a decent snack situation --- you bring yourself.
